%=========================
% Sections/01-GioiThieu.tex
% Chương 1 – Giới thiệu
%=========================

\section{Giới thiệu}

\subsection{Động lực}
\indent Luyện tập IELTS Speaking đòi hỏi người học nhận được phản hồi về phát âm, độ trôi chảy, từ vựng và ngữ pháp, đồng thời theo dõi sự tiến bộ qua nhiều bài nói. Tuy nhiên, việc luyện tập cùng giáo viên phụ thuộc vào thời gian và điều kiện tiếp cận của mỗi học viên. Các công cụ đánh giá tự động có thể hỗ trợ luyện tập thường xuyên, nhưng kết quả còn chịu ảnh hưởng của chất lượng audio, nhận dạng tiếng nói và khả năng đánh giá của mô hình. Bên cạnh đó, điểm số riêng lẻ chưa đủ để xác định nội dung luyện tập phù hợp với từng người học.

\indent Từ nhu cầu đó, đề tài xây dựng hệ thống AI hỗ trợ chấm điểm và cá nhân hóa luyện tập kỹ năng Nói IELTS, kết hợp mô hình học sâu chuyên biệt cho phát âm với mô hình ngôn ngữ lớn để đánh giá các tiêu chí còn lại. Kết quả được tổng hợp thành hồ sơ năng lực theo thời gian, làm cơ sở đề xuất nội dung luyện tập theo điểm yếu và tiến độ. Hệ thống hướng đến tích hợp thử nghiệm vào nền tảng WISPACE và sử dụng AIHub để hỗ trợ cung cấp, quản lý dịch vụ AI qua API.

\subsection{Mục tiêu}
\indent Mục tiêu của đề tài là nghiên cứu và xây dựng hệ thống hỗ trợ đánh giá IELTS Speaking, cung cấp phản hồi và cá nhân hóa quá trình luyện tập. Trong giai đoạn Đồ án chuyên ngành, đề tài tập trung khảo sát, phân tích, thiết kế và hiện thực các chức năng, giải pháp được lựa chọn. Các mục tiêu chính gồm:
\begin{itemize}
    \item Kết hợp mô hình đánh giá phát âm với LLM và bằng chứng âm học để hỗ trợ chấm bốn tiêu chí: Pronunciation, Fluency \& Coherence, Lexical Resource và Grammatical Range \& Accuracy.
    \item Cung cấp điểm, nhận xét và dẫn chứng; lưu lịch sử để hình thành hồ sơ năng lực và đề xuất luyện tập cá nhân hóa.
    \item Xây dựng API và các chức năng cần thiết cho việc tích hợp vào Wispace Speaking; quản lý dịch vụ AI và dữ liệu học thuật thông qua AIHub.
    \item Hiện thực, kiểm thử các chức năng và so sánh các giải pháp trong phạm vi giai đoạn 1; đánh giá kết quả và trình bày rõ những giới hạn còn tồn tại.
\end{itemize}

\subsection{Phạm vi}

\indent Đề tài tập trung vào bài toán chấm điểm và cá nhân hóa luyện tập IELTS Speaking. Wispace Speaking cung cấp môi trường luyện tập, thi thử và theo dõi kết quả; AIHub là hệ thống độc lập phục vụ quản lý API mô hình cho đối tác B2B và công cụ học thuật nội bộ. Phạm vi báo cáo tương ứng với giai đoạn 1, làm cơ sở hoàn thiện hệ thống trong giai đoạn tiếp theo.

\begin{itemize}
    \item \textbf{Trong phạm vi:}
    \begin{itemize}
        \item Tìm hiểu yêu cầu nghiệp vụ, kiến thức nền tảng, bài báo và hệ thống liên quan; phân tích ưu điểm, hạn chế và khoảng trống cần giải quyết.
        \item Phân tích, thiết kế kiến trúc, dữ liệu, pipeline chấm Speaking và cơ chế cá nhân hóa dựa trên lịch sử học tập.
        \item Hiện thực và kiểm thử ít nhất $2n$ chức năng với $n$ là số thành viên; với nhóm ba thành viên, số chức năng tối thiểu là sáu. Hiện thực và đánh giá ít nhất hai giải pháp cho bài toán mục tiêu.
        \item Xây dựng API, kiểm thử tích hợp và thực hiện thử nghiệm trên WISPACE theo điều kiện dữ liệu, quyền truy cập và môi trường được cấp.
        \item Đánh giá phần đã hiện thực, phân biệt mục tiêu thiết kế với kết quả thực nghiệm và xác định hướng hoàn thiện ở giai đoạn 2.
    \end{itemize}

    \item \textbf{Ngoài phạm vi:}
    \begin{itemize}
        \item Cấp chứng nhận hoặc thay thế kết quả thi IELTS chính thức.
        \item Chấm toàn bộ bốn kỹ năng IELTS; dữ liệu Writing trong AIHub chỉ phục vụ nhóm chức năng học thuật liên quan.
        \item Khẳng định hiệu quả cải thiện dài hạn hoặc khả năng vận hành tải lớn khi chưa có thử nghiệm tương ứng.
        \item Hoàn thiện toàn bộ chức năng và triển khai thương mại đầy đủ ngay trong giai đoạn 1.
    \end{itemize}

    \item \textbf{Giả định:}
    \begin{itemize}
        \item Người học sử dụng thiết bị có micro và trình duyệt hỗ trợ ghi âm, với kết nối Internet phù hợp để nộp bài.
        \item Dữ liệu câu hỏi, audio và điểm đối chiếu được sử dụng trong phạm vi cho phép; việc khai thác dữ liệu cá nhân phục vụ huấn luyện có sự đồng ý tương ứng.
        \item Nền tảng WISPACE cung cấp điều kiện tích hợp và thử nghiệm theo phạm vi được thống nhất.
        \item Mô hình và dịch vụ AI được triển khai trong môi trường thực nghiệm; lỗi nhận dạng và sai lệch chấm điểm được kiểm tra, không mặc định đầu ra AI là chính xác.
    \end{itemize}
\end{itemize}

\subsection{Ý nghĩa của đề tài}
\subsubsection{Ý nghĩa về mặt thực tiễn}
\indent Đề tài hướng đến hỗ trợ người học luyện nói thường xuyên, tiếp cận phản hồi và theo dõi tiến bộ trên cùng nền tảng. Hồ sơ năng lực giúp đề xuất nội dung phù hợp với điểm yếu và mục tiêu của từng học viên. Việc đóng gói thành API tạo điều kiện tích hợp vào Wispace và quản lý cung cấp dịch vụ thông qua AIHub.

\subsubsection{Ý nghĩa về mặt khoa học}
\indent Đề tài nghiên cứu việc kết hợp mô hình âm học với LLM để đánh giá Speaking và sử dụng kết quả cho cá nhân hóa luyện tập. Quá trình đối chiếu các giải pháp giúp phân tích ảnh hưởng của transcript ASR, đặc trưng âm học và phương pháp tổng hợp điểm. Kết quả được trình bày cùng điều kiện đánh giá, giới hạn dữ liệu và phạm vi kết luận, bảo đảm minh bạch khi sử dụng AI.

\subsection{Cấu trúc báo cáo}

\indent Báo cáo được tổ chức thành các chương chính như sau:
\begin{itemize}
    \item Chương 1: Giới thiệu động lực, mục tiêu, phạm vi và ý nghĩa của đề tài.
    \item Chương 2: Trình bày kiến thức nền tảng về IELTS Speaking, xử lý tiếng nói, đánh giá tự động và dữ liệu phục vụ AI.
    \item Chương 3: Khảo sát công trình, hệ thống và công nghệ liên quan; phân tích ưu điểm, hạn chế và cơ sở đề xuất giải pháp.
    \item Chương 4: Giới thiệu công nghệ được lựa chọn và vai trò trong Wispace, AIHub.
    \item Chương 5: Khai thác yêu cầu chức năng, phi chức năng và dữ liệu của hai hệ thống.
    \item Chương 6: Phân tích tác nhân, ca sử dụng, luồng hoạt động và luồng dữ liệu.
    \item Chương 7: Trình bày thiết kế kiến trúc, dữ liệu và tích hợp các thành phần AI.
    \item Chương 8: Mô tả phần cài đặt, hiện thực và kiểm thử trong phạm vi giai đoạn 1.
    \item Các chương cuối: Đánh giá kết quả đạt được, tổng kết hạn chế và đề xuất hướng hoàn thiện ở giai đoạn tiếp theo.
\end{itemize}
